%!TEX program = lualatex
%!TEX encoding = UTF-8 Unicode  
%!TEX root = chapter04.tex  
\documentclass[main.tex]{subfiles} 
\begin{document} 
%----------------------------------------------------------------------------------------
%	 
%----------------------------------------------------------------------------------------
\newpage 
\loadgeometry{marginlayout}  
\pagestyle{scrheadings} % Print headers again  
\KOMAoptions{
	headwidth=textwithmarginpar,
	headsepline=true,
	headsepline= 0.5pt,
	footwidth=textwithmarginpar,
}  
\onehalfspacing 

%----------------------------------------------------------------------------------------
%	CHAPTER 4
%----------------------------------------------------------------------------------------
\setcounter{chapter}{3}

\chapter{Inéquations}
\section{Vocabulaire}
 Une \textbf{inéquation à une inconnue} est une inégalité dans laquelle apparaît une lettre.
 
 \textbf{Une solution} de l’inéquation est une valeur de la ou les inconnues pour lesquelles l’inégalité est vraie. 
 \begin{example}%\hfill \\
 	Soit l'inéquation $4x+7<x^2$ d'inconnue $x$.
 	\begin{enumerate}[label=\alph*),leftmargin=*]
 		\item  		$x=0$ n'est pas solution de l'équation car l'égalité $4\times 0+7<0^2$ est fausse.
 		\item $x=6$ est une solution de l'équation, car $4\times 6+7<6^2$ est vraie.
 	\end{enumerate}
 \end{example}
 \begin{example}%\hfill \\
	Soit l'inéquation $7x-12\geqslant x^2$ d'inconnue $x$.
	\begin{enumerate}[label=\alph*),leftmargin=*]
		\item $x=10$ n'est pas solution de l'équation car l'inégalité $7\times 10-12\geqslant 10^2$ est fausse.
		\item $x=3$,  $x=\numprint{3.5}$  et $x=4$ sont solutions de l'inéquation, car $7\times 3-12\geqslant3^2$ est vraie. 
	\end{enumerate}
\end{example}
 \begin{definition} Résoudre une équation dans $\mathbb{R}$ c'est trouver toutes les valeurs  réelles des  inconnues qui rendent l'inégalité vraie. 
 \end{definition} 
\begin{definition}
	Deux inéquations sont dites \textbf{équivalentes} (symbole $\iff$) si elles ont le même ensemble de solutions c.à.d elles sont vraies pour les mêmes valeurs de $x$.
\end{definition} 
\begin{example}\hfill 
	\begin{enumerate}[label=\alph*),leftmargin=*]
		\item Les inéquations $2x>4$ et $2x-4>0$ d'inconnue $x$ sont équivalentes.
		\item  $x=-2$, $x=-3$, $x=-4$... ne sont pas solutions de $x^2\leqslant4$.\\
		Les inéquations $x\leqslant 2$ et $x^2\leqslant 4$ ne sont pas équivalentes.
	\end{enumerate}
\end{example}
\vfill 

 

%----------------------------------------------------------------------------------------
%	EXERCICES
%----------------------------------------------------------------------------------------

\newpage 
\loadgeometry{widelayout}  
\pagestyle{scrheadings} % Print headers again  
\KOMAoptions{
	headwidth=textwithmarginpar,
	headsepline=true,
	headsepline= 0.5pt,
	footwidth=textwithmarginpar,
} 
\onehalfspacing
\subsection{Exercices : mise en inéquations}
\begin{example}Traduire le problème posé par une inéquation.\smallskip
	
\noindent \parbox{0.48\linewidth}{	\smallskip
La loi impose d'avoir au minimum 12\si{\meter^2} d'espace par cochon dans un enclos. Combien de cochons peut accueillir un enclos de 108~\si{\meter^2} ?\\
$x = \ldots$\\
On chercher $x\in \ldots$ qui vérifie :\\ 
	$\begin{WithArrows} 
		&  \rule{0pt}{1.8em}   \\
		&  \rule{0pt}{1.8em}   \\
		&  \rule{0pt}{1.8em}   \\
		&  \rule{0pt}{1.8em}   \\
		&  \rule{0pt}{1.8em}   \\
		&  \rule{0pt}{1.1em}  
	\end{WithArrows}$ 	
}\hfill \parbox{0.48\linewidth}{\smallskip	
Pizza TropBien vend sa pizza \EUR{5} l'unité. La fabrication d'une pizza lui revient à  \EUR{2.3} diminué de \EUR{0.02} par pizza vendue. Combien doit-il vendre de pizza pour faire au delà de \EUR{180} de profits ?\\
$x = \ldots$\\
On chercher $x\in \ldots$ qui vérifie :\\  
	$\begin{WithArrows} 
		&  \rule{0pt}{1.8em}   \\
		&  \rule{0pt}{1.8em}   \\
		&  \rule{0pt}{1.8em}   \\
		&  \rule{0pt}{1.8em}    \\
		&  \rule{0pt}{1.8em}   
	\end{WithArrows}$
} 
\end{example} 
\begin{exercise} \label{chapter08:exo:vocabulaire01}
	Même consignes  
	\begin{enumerate}[label=\alph*), leftmargin=*]
%		\item La loi impose d'avoir au minimum 12\si{\meter^2} d'espace par cochon dans un enclos. Combien de cochons peut accueillir un enclos de 108~\si{\meter^2} ?
		\item Pour valider le module de Pix, un élève doit accumuler 120 points en 2 tests. Helga a 54 points au premier test mais ne valide pas le module. Quel est le plus grand score possible au second test ?
		\item Harold veut acheter un vélo à \EUR{310}. Il a \EUR{65} et met de côté \EUR{45} chaque semaine grâce à son job d'été. Combien de semaines avant de pouvoir acheter ce vélo ?
		\item Eugene a \EUR{50} et met de côté \EUR{6} par semaine. Lila n'a pas d'économie et met de côté \EUR{9} par semaine. Combien de semaines sont nécessaires pour que Lila ait plus d'argent que Kyle. 
		\item Arnold a \EUR{18}. Il veut acheter des cupcakes à \EUR{1.5} pièce. Quel est le nombre maximal de cupcakes qu'il peut s'offrir ?
		\item Un taxi prend \EUR{5} de frais de service et \EUR{3} par \si{\kilo\meter} de trajet. Quelle est le plus long trajet que peut se payer Rhonda avec \EUR{71} ?
%		\item Pizza Trop Bien vend sa pizza à \EUR{5} l'unité. La fabrication lui coûte \EUR{1.5} par pizza. Combien doit-il vendre de pizza pour faire au moins \EUR{180} de profits ?
	\end{enumerate}
\end{exercise} 
\begin{exercise}[Vérifier si une valeur est solution d'une inéquation à 1 inconnue]\label{chapter08:exo:vocabulaire02}\hfill 
	
	\noindent\QCM[VF,Alterne,Stretch=1.0, Largeur=0.75cm]{%
		%		{$70$ est solution de l'équation $x^{300}=1$ d'inconnue $x$}&2,%
		{$3$ est une solution de l'inéquation $2x+1<5$ d'inconnue $x$}&2,%
		{$2$ est une solution de l'inéquation $2x+1< 5$ d'inconnue $x$}&2,%
		{$-5$ est une solution de l'inéquation $7-x > x^2-13$ d'inconnue $x$}&1,%
		{$4$ est une solution de l'inéquation $x\leqslant 4$ d'inconnue $x$}&1,% 
		{Les inéquations $x-3>0$ et $x>3$ sont équivalentes.}&1,% 
		{Les inéquations $3x \leqslant 1$ et $x\leqslant-2$ d'inconnue $x$ sont équivalentes.}&2,% 
		{Les inéquations $3x \leqslant 0$ et $x\leqslant -3$ d'inconnue $x$ sont équivalentes.}&2,%   
	}  
\end{exercise}

%----------------------------------------------------------------------------------------
%	 
%----------------------------------------------------------------------------------------
\newpage
\begin{proof}[solution de l'exercice \ref{chapter08:exo:vocabulaire01}]\hfill 
	\begin{enumerate}[label=\alph*), leftmargin=*] 
	\item $x=$ note au test, $x\in\mathbb{N}$. $x$ vérifie  $54+x<120$. \xout{$54+x\leqslant 120$} 
	\item $x=$ nombre de semaines, $x\in\mathbb{N}$. $x$ vérifie  $65+45x\geqslant 310$. \xout{$65+45x > 310$}
	\item   $x=$ nombre de semaines, $x\in\mathbb{N}$. $x$ vérifie $50+6x<9x$. \xout{$50+6x\geqslant 9x$} 
	\item $x=$ nombre de cupcakes,  $x\in\mathbb{N}$. $x$ vérifie $18\geqslant 1.5x$. \xout{$18> 1.5x$} 
	\item $x=$ longueur du trajet en km,  $x\in\mathbb{R}$. $x$ vérifie $71\geqslant 3x+5$. \xout{$71> 3x+5$}  
\end{enumerate}
\end{proof}
\begin{proof}[solution de l'exercice \ref{chapter08:exo:vocabulaire02}]\hfill 
	
	\noindent\QCM[VF,Alterne,Stretch=1.0,Solution, Largeur=0.75cm]{%
		%		{$70$ est solution de l'équation $x^{300}=1$ d'inconnue $x$}&2,%
		{$3$ est une solution de l'inéquation $2x+1<5$ d'inconnue $x$}&2,%
		{$2$ est une solution de l'inéquation $2x+1< 5$ d'inconnue $x$}&2,%
		{$-5$ est une solution de l'inéquation $7-x > x^2-13$ d'inconnue $x$}&1,%
		{$4$ est une solution de l'inéquation $x\leqslant 4$ d'inconnue $x$}&1,% 
		{Les inéquations $x-3>0$ et $x>3$ sont équivalentes.}&1,% 
		{Les inéquations $3x \leqslant 1$ et $x\leqslant-2$ d'inconnue $x$ sont équivalentes.}&2,% 
		{Les inéquations $3x \leqslant 0$ et $x\leqslant -3$ d'inconnue $x$ sont équivalentes.}&2,%   
	}   
\end{proof}
%----------------------------------------------------------------------------------------
%	COURS
%----------------------------------------------------------------------------------------

\newpage 
\loadgeometry{marginlayout}  
\pagestyle{scrheadings} % Print headers again  
\KOMAoptions{
	headwidth=textwithmarginpar,
	headsepline=true,
	headsepline= 0.5pt,
	footwidth=textwithmarginpar,
} 
\onehalfspacing
\section{Intervalles}
 

\begin{example} Soit l'inéquation  $x\leqslant 4$ d'inconnue $x$.
	L'ensemble des solutions dans $\mathbb{R}$ est  $ I = \{ x \in \mathbb{R} \quad |     x \leqslant 4 \} $.
	\begin{marginfigure}
		\begin{tikzpicture}[scale=0.6]
			\tkzSetUpPoint[size=2, fill=black]
			\tkzfctset{tan style/.style={->,>=latex, color=red, line width=1.2pt}} 
			\tkzInit[xmin=-1.5, xmax=6.0, xstep=1, ymax=0.75, ymin=-0.95]
			\tkzDrawX[label={$x$},above=5pt , line width=1.2pt] %  noticks, 
			\tkzClip
			%		\tkzHTick[mark=text,mark options={color=black}, text mark={\LARGE ]}]{4.01}
			%		\tkzHTick[mark=text,mark options={color=black}, text mark={ {\LARGE [}}]{3.01}
			\tkzDefPoint(0,0){O}				\tkzLabelPoint[below=4pt](O){\scriptsize$0$}
			\tkzDefPoint(-2,0){A}
			\tkzDefPoint(4.0,0){B}
			\tkzDrawSegments[ -*, line width=3 pt, color=black](A,B)
			%		\tkzLabelPoint[below=4pt](A){\scriptsize$3$}
			\tkzLabelPoint[below=4pt](B){\scriptsize$4$} 
		\end{tikzpicture} 
		\caption{$I= \interval[open left]{-\infty}{4}$}
		\label{fig:chap03intervalleA}
	\end{marginfigure}
	On le notera $I= \interval[open left]{-\infty}{4}$. Intervalle de $-\infty$ à $4$ fermé en $4$.
\end{example}
\begin{example} 
	La double inéquation $3 \leqslant x \leqslant 5$
	d'inconnue $x$ a pour solution $J = \{ x \in \mathbb{R} \quad |  \quad 3 \leqslant x \leqslant 5 \}$.

	se lit :\og l'ensemble des $x$ dans $\mathbb{R}$ avec $3 \leqslant x \leqslant 5$ \fg{}.
	\begin{marginfigure} 
		\begin{tikzpicture}[scale=0.6]
			\tkzSetUpPoint[size=2, fill=black]
			\tkzfctset{tan style/.style={->,>=latex, color=red, line width=1.2pt}} 
			\tkzInit[xmin=-1.5, xmax=6.0, xstep=1, ymax=0.75, ymin=-0.95]
			\tkzDrawX[label={$x$},above=5pt , line width=1.2pt] %  noticks, 
			\tkzClip
			%			\tkzHTick[mark=text,mark options={color=black}, text mark={\LARGE ]}]{5.01}
			%			\tkzHTick[mark=text,mark options={color=black}, text mark={ {\LARGE [}}]{3.01}
			\tkzDefPoint(0,0){O}				\tkzLabelPoint[below=4pt](O){\scriptsize$0$}
			\tkzDefPoint(3,0){A}
			\tkzDefPoint(5.0,0){B}
			\tkzDrawSegments[ *-*, line width=3 pt, color=black](A,B)
			\tkzLabelPoint[below=4pt](A){\scriptsize$3$}
			\tkzLabelPoint[below=4pt](B){\scriptsize$5$} 
		\end{tikzpicture} 
		\caption{$J=\interval{3}{5}$}
		\label{fig:chap03intervalleB}
	\end{marginfigure}
	On le notera $J=\interval{3}{5}$. 
\end{example}



\begin{table*}[!hb]
	\begin{tabularx}{1\columnwidth}{|c?c?>{}X|}\hline
		Intervalle & Inégalité & Représentation sur droite graduée\centering   \tabularnewline  \hline \hline 
		$x\in [a;b]$ & $a\leqslant x \leqslant b$ &
		\centering  \begin{tikzpicture} 
			\tkzInit[xmin=-3.5, xmax=6.0, xstep=1, ymax=0.75, ymin=-0.75]
			\tkzDrawX[  label={$x$},above=5pt ] %noticks, 
			\tkzClip
			%			\tkzHTick[mark=text,mark options={color=black}, text mark={\LARGE ]}]{4.01}
			%			\tkzHTick[mark=text,mark options={color=black}, text mark={ {\LARGE [}}]{-2.01}
			
			\tkzDefPoint(-2,0){A}
			\tkzDefPoint(4.0,0){B}
			\tkzDrawSegments[ *-*,  line width=3 pt, color=black](A,B)
			\tkzLabelPoint[below=8pt](A){$a$}
			\tkzLabelPoint[below=8pt](B){$b$}
			\tkzLabelPoint[above=10pt](B){}
		\end{tikzpicture}     \tabularnewline \hline
		%
		$x\in ]a;b[$ & $a < x < b$ &
		\begin{tikzpicture} 
			\tkzInit[xmin=-3.5, xmax=6.0, xstep=1, ymax=0.75, ymin=-0.75]
			\tkzDrawX[  label={ },above=5pt ] %noticks, 
			\tkzClip
			%			\tkzHTick[mark=text,mark options={color=black}, text mark={\LARGE [}]{4.01}
			%			\tkzHTick[mark=text,mark options={color=black}, text mark={ {\LARGE ]}}]{-2.01}
			
			\tkzDefPoint(-2,0){A}
			\tkzDefPoint(4.0,0){B}
			\tkzDrawSegments[ o-o,  line width=3 pt, color=black](A,B)
			\tkzLabelPoint[below=8pt](A){$a$}
			\tkzLabelPoint[below=8pt](B){$b$}
			\tkzLabelPoint[above=10pt](B){}
		\end{tikzpicture} \centering   \tabularnewline \hline
		%
		$x\in [a;b[$ & $a\leqslant x < b$ &
		\begin{tikzpicture} 
			\tkzInit[xmin=-3.5, xmax=6.0, xstep=1, ymax=0.75, ymin=-0.75]
			\tkzDrawX[  label={ },above=5pt ] %noticks, 
			\tkzClip  
			%			\tkzHTick[mark=text,mark options={color=black}, text mark={\LARGE [}]{4.01}
			%			\tkzHTick[mark=text,mark options={color=black}, text mark={ {\LARGE [}}]{-2.01}
			
			\tkzDefPoint(-2,0){A}
			\tkzDefPoint(4.0,0){B}
			\tkzDrawSegments[  *-o, line width=3 pt, color=black](A,B)
			\tkzLabelPoint[below=8pt](A){$a$}
			\tkzLabelPoint[below=8pt](B){$b$}
			\tkzLabelPoint[above=10pt](B){}
		\end{tikzpicture} \centering   \tabularnewline \hline
		%
		$x\in ]a;b]$ & $a< x \leqslant b$ &
		\begin{tikzpicture} 
			\tkzInit[xmin=-3.5, xmax=6.0, xstep=1, ymax=0.75, ymin=-0.75]
			\tkzDrawX[  label={ },above=5pt ] %noticks, 
			\tkzClip
			%			\tkzHTick[mark=text,mark options={color=black}, text mark={\LARGE ]}]{4.01}
			%			\tkzHTick[mark=text,mark options={color=black}, text mark={ {\LARGE ]}}]{-2.01}
			
			\tkzDefPoint(-2,0){A}
			\tkzDefPoint(4.0,0){B}
			\tkzDrawSegments[ o-*, line width=3 pt, color=black](A,B)
			\tkzLabelPoint[below=8pt](A){$a$}
			\tkzLabelPoint[below=8pt](B){$b$}
			\tkzLabelPoint[above=10pt](B){}
		\end{tikzpicture} \centering   \tabularnewline \hline
	\end{tabularx} 
	\caption{Les différentes variantes d'intervalles bornés par $a$ et $b\in\mathbb{R}$} 
\end{table*}

\begin{table*}[!hb]%[tbp]
	\begin{tabularx}{1\columnwidth}{|c?c?>{}X|}\hline
		Intervalle & Inégalité & Représentation sur droite graduée\centering   \tabularnewline  \hline \hline 
		$x\in [a;+\infty[$ & $x\geqslant a$ &
		\begin{tikzpicture} 
			\tkzInit[xmin=-3.5, xmax=6.0, xstep=1, ymax=0.75, ymin=-0.75]
			\tkzDrawX[  label={ },above=5pt ] %noticks, 
			%\tkzHTick[mark=text,mark options={color=black}, text mark={\LARGE ]}]{4.01}
			%			\tkzHTick[mark=text,mark options={color=black}, text mark={ {\LARGE [}}]{-2.01}
			
			\tkzDefPoint(-2,0){A}
			\tkzDefPoint(6.35,0){B}
			\tkzDrawSegments[ *-,  line width=3 pt, color=black](A,B)
			\tkzLabelPoint[below=8pt](A){$a$}
			%\tkzLabelPoint[below=8pt](B){$b$}
			\tkzLabelPoint[above=10pt](B){}
			\tkzClip 
		\end{tikzpicture} \centering   \tabularnewline \hline
		%
		$x\in ]a;+\infty[$ & $x > a$ &
		\begin{tikzpicture} 
			\tkzInit[xmin=-3.5, xmax=6.0, xstep=1, ymax=0.75, ymin=-0.75]
			\tkzDrawX[  label={ },above=5pt ] %noticks, 
			%\tkzHTick[mark=text,mark options={color=black}, text mark={\LARGE ]}]{4.01}
			%			\tkzHTick[mark=text,mark options={color=black}, text mark={ {\LARGE ]}}]{-2.01}
			
			\tkzDefPoint(-2,0){A}
			\tkzDefPoint(6.35,0){B}
			\tkzDrawSegments[ o- ,  line width=3 pt, color=black](A,B)
			\tkzLabelPoint[below=8pt](A){$a$}
			%\tkzLabelPoint[below=8pt](B){$b$}
			\tkzLabelPoint[above=10pt](B){}
			\tkzClip
		\end{tikzpicture} \centering   \tabularnewline \hline\hline
		%
		%
		$x\in ]-\infty; b]$ & $x \leqslant b$ &
		\begin{tikzpicture} 
			\tkzInit[xmin=-3.5, xmax=6.0, xstep=1, ymax=0.75, ymin=-0.75]
			\tkzDrawX[  label={ },above=5pt ] %noticks, 
			\tkzClip
			%			\tkzHTick[mark=text,mark options={color=black}, text mark={\LARGE ]}]{4.01}
			%\tkzHTick[mark=text,mark options={color=black}, text mark={ {\LARGE [}}]{-2.01}		     
			\tkzDefPoint(-3.5,0){A}
			\tkzDefPoint(4.0,0){B}
			\tkzDrawSegments[   -*, line width=3 pt, color=black](A,B)
			%\tkzLabelPoint[below=8pt](A){$a$}
			\tkzLabelPoint[below=8pt](B){$b$}
			\tkzLabelPoint[above=10pt](B){}
		\end{tikzpicture} \centering   \tabularnewline \hline
		%
		$x\in \interval[open,scaled]{-\infty}{b}$ & $x< b$ &
		\begin{tikzpicture} 
			\tkzInit[xmin=-3.5, xmax=6.0, xstep=1, ymax=0.75, ymin=-0.75]
			\tkzDrawX[  label={ },above=5pt ] %noticks, 
			\tkzClip
			%			\tkzHTick[mark=text,mark options={color=black}, text mark={\LARGE [}]{4.01}
			%\tkzHTick[mark=text,mark options={color=black}, text mark={ {\LARGE ]}}]{-2.01}
			\tkzDefPoint(-3.5,0){A}
			\tkzDefPoint(4.0,0){B}
			\tkzDrawSegments[ -o,  line width=3 pt, color=black](A,B)
			%\tkzLabelPoint[below=8pt](A){$a$}
			\tkzLabelPoint[below=8pt](B){$b$}
			\tkzLabelPoint[above=10pt](B){}
		\end{tikzpicture} \centering   \tabularnewline \hline
		%
	\end{tabularx} 
	\caption{Les différentes variantes d'intervalles infinis} 
\end{table*}
 
%----------------------------------------------------------------------------------------
%	EXERCICES
%----------------------------------------------------------------------------------------

\newpage 
\loadgeometry{widelayout}  
\pagestyle{scrheadings} % Print headers again  
\KOMAoptions{
	headwidth=textwithmarginpar,
	headsepline=true,
	headsepline= 0.5pt,
	footwidth=textwithmarginpar,
} 
\onehalfspacing
\subsection{Exercices : Intervalles}
\begin{exercise} \hfill 
	
	Sans calculatrice, compléter par l'un des symboles : $\in$, $\notin$.
	
	$-\numprint{3.1}~\dots~ \interval[scaled]{-4}{-3}$ ;\hfill $\numprint{2.3}\times10^{-2} ~\dots~\interval[scaled]{2}{3}$;\hfill $\frac{1}{4} ~\dots~\interval[scaled]{\frac{1}{3}}{\frac{1}{2}}$;\hfill $-\frac{4}{5} ~\dots~\interval[scaled]{-1}{-\frac{3}{4}}$;\hfill
\end{exercise} 
\begin{exercise} \hfill 
	
	Sans calculatrice, compléter par  $\in$, $\notin$ et $\subset$ ou $\supset$.
	
	$\frac{17}{4}~\dots~ \interval[open, scaled]{4}{5}$ ;\hfill
	$\numprint{0.333}~\dots~ \interval[ open right, scaled]{\frac{1}{3}}{1}$ ;\hfill
	$\mathbb{Z}~\dots~\mathbb{R}$;\hfill 
	$\dfrac{2}{3}~\dots~  \mathbb{D}$ \hfill
	
	$\sqrt{8}~\dots~ \interval[open, scaled]{2}{3}$ ;\hfill
	$\frac{3}{8}~\dots~ \interval[ scaled]{\frac{3}{9}}{\frac{3}{7}}$ ;\hfill 
	$\frac{11}{3} ~\dots~  \interval[ open left, scaled]{\frac{21}{6}}{5}$ \hfill
	$  \interval[ open left, scaled]{\frac{21}{6}}{5} ~\dots~ \interval[ scaled]{\frac{11}{3}}{4}  $ \hfill
\end{exercise}


\begin{exercise}  
	Compléter le tableau suivant :\\
	\begin{tabularx}{1\columnwidth}{|c?c?X?>{}p{6cm}|}\hline
		Intervalle & Inégalité(s) & Représentation sur droite réelle\centering   &  Phrase\centering\tabularnewline  \hline \hline 
		$x\in [-3;5]$ &  &
		\begin{tikzpicture}[xscale=0.45]
			\tkzInit[xmin=-4.5, xmax=6.0, xstep=1, ymax=0.75, ymin=-0.73]
			\tkzDrawX[ noticks, label={ },above=5pt ] %noticks, 
			\tkzClip
			%			\tkzHTick[mark=text,mark options={color=black}, text mark={\LARGE ]}]{5.01}
			%			\tkzHTick[mark=text,mark options={color=black}, text mark={ {\LARGE [}}]{-3.01}
			
			\tkzDefPoint(-3,0){A}
			\tkzDefPoint(5.0,0){B}
			\tkzDrawSegments[ *-*, line width=3 pt, color=black](A,B)
			\tkzLabelPoint[below=8pt](A){\scriptsize$-3$}
			\tkzLabelPoint[below=8pt](B){\scriptsize$5$}
			\tkzLabelPoint[below=2pt](0,0){$0$}
		\end{tikzpicture}   &   \tabularnewline \hline
		%
		& $  x<3$ 
		& 	\centering  \begin{tikzpicture}[xscale=0.5]
			\tkzInit[xmin=-3.5, xmax=6.0, xstep=1, ymax=0.75, ymin=-0.73]
			\tkzDrawX[ noticks, label={ },above=5pt ] %noticks, 
			\tkzClip
		\end{tikzpicture}
		& \\ \hline
		& 
		&	\centering  \begin{tikzpicture}[xscale=0.5]
			\tkzInit[xmin=-3.5, xmax=6.0, xstep=1, ymax=0.75, ymin=-0.73]
			\tkzDrawX[ noticks, label={ },above=5pt ] %noticks, 
			\tkzClip
		\end{tikzpicture}
		& Intervalle de 4 à 6, fermé en 4 et ouvert en 6. \centering\tabularnewline \hline
		$[2;+\infty[$ &  &
		\begin{tikzpicture}[xscale=0.5]
			\tkzInit[xmin=-3.5, xmax=6.0, xstep=1, ymax=0.75, ymin=-0.73]
			\tkzDrawX[ noticks, label={ },above=5pt ] %noticks, 
			\tkzClip
		\end{tikzpicture}   &   \tabularnewline \hline
		& $-3< x \leqslant -1$ &
		\begin{tikzpicture}[xscale=0.5]
			\tkzInit[xmin=-3.5, xmax=6.0, xstep=1, ymax=0.75, ymin=-0.73]
			\tkzDrawX[ noticks, label={ },above=5pt ] %noticks, 
			\tkzClip
		\end{tikzpicture}   &   \tabularnewline \hline
		& 
		&	  \begin{tikzpicture}[xscale=0.5]
			\tkzInit[xmin=-3.5, xmax=6.0, xstep=1, ymax=0.75, ymin=-0.73]
			\tkzDrawX[ noticks, label={ },above=5pt ] %noticks, 
			\tkzClip
		\end{tikzpicture}
		& Intervalle de $-\infty$ à 5, fermé en 5. \centering\tabularnewline \hline
		& $-3\leqslant x \leqslant -1$ &
		\begin{tikzpicture}[xscale=0.5]
			\tkzInit[xmin=-3.5, xmax=6.0, xstep=1, ymax=0.75, ymin=-0.73]
			\tkzDrawX[ noticks, label={ },above=5pt ] %noticks, 
			\tkzClip
		\end{tikzpicture}   &   \tabularnewline \hline
		& $5 \geqslant  x >1 $ &
		\begin{tikzpicture}[xscale=0.5]
			\tkzInit[xmin=-3.5, xmax=6.0, xstep=1, ymax=0.75, ymin=-0.73]
			\tkzDrawX[ noticks, label={ },above=5pt ] %noticks, 
			\tkzClip
		\end{tikzpicture}   &   \tabularnewline \hline
		& $x\geqslant -\frac{3}{4}$ &
		\begin{tikzpicture}[xscale=0.5]
			\tkzInit[xmin=-3.5, xmax=6.0, xstep=1, ymax=0.75, ymin=-0.73]
			\tkzDrawX[ noticks, label={ },above=5pt ] %noticks, 
			\tkzClip
		\end{tikzpicture}   &   \tabularnewline \hline
		& $-4 > x > -7 $ &
		\begin{tikzpicture}[xscale=0.5]
			\tkzInit[xmin=-3.5, xmax=6.0, xstep=1, ymax=0.75, ymin=-0.73]
			\tkzDrawX[ noticks, label={ },above=5pt ] %noticks, 
			\tkzClip
		\end{tikzpicture}   &   \tabularnewline \hline
	\end{tabularx}
\end{exercise} 
 
\newpage 
\begin{example}[Intersection d'intervalles] est toujours un intervalle. L'union de deux intervalles d'intersection vide n'est pas un intervalle.\hfill
	
\noindent
\parbox{0.4\linewidth}{
\noindent $A=\interval[open left]{-\infty}{1}$ et $B=\interval[open]{0}{\infty}$.

$A\cap B=\interval[open left]{0}{1}$ \\
$A\cup B=\interval[open]{-\infty}{\infty}$ 
}
\hfill 
\parbox{0.6\linewidth}{
\begin{tikzpicture}[xscale=0.8]
	\tkzInit[xmin=-6, xmax=6.0, xstep=1, ymax=1.75, ymin=-0.73]
	\tkzDrawX[ noticks, label={ },above=5pt, line width=2pt ] %noticks, 
	\tkzClip
	\tkzDefPoint(-6.0,0){A1}
	\tkzDefPoint(2.0,0){A2}
	\tkzDefPoint(-2.0,0){B1}
	\tkzDefPoint(6.0,0){B2}
	\draw    (-6,1) to[out=0,in=150]  node[pos=0.5, fill=white ] {$A$} (A2)  ;
	\draw    (B1) to[out=30,in=180]  node[pos=0.5, fill=white ] {$B$} (6,1)  ;
	\tkzDrawPoint[shape=circle,size=8](A2);
%	\tkzDrawPoint[shape=circle,fill=white,size=8](A2);
%	\tkzLabelPoint[below](A1){$0$}
	\tkzLabelPoint[below](A2){$1$}
	\tkzDrawPoint[shape=circle,fill=white, size=8](B1);
%	\tkzDrawPoint[shape=circle,  size=8](B2);
	\tkzLabelPoint[below](B1){$0$}
%	\tkzLabelPoint[below](B2){$1$}
\end{tikzpicture} 	
}

\vspace{-0.75em}
\noindent
\parbox{0.4\linewidth}{
	\noindent $A=\interval[open]{-3}{2}$ et $B=\interval[]{1}{5}$.
	
	$A\cap B=\interval[open right]{1}{2}$ \\
	$A\cup B=\interval[open left]{-3}{5}$ 
}
\hfill 
\parbox{0.6\linewidth}{
	\begin{tikzpicture}[xscale=0.8]
		\tkzInit[xmin=-6, xmax=6.0, xstep=1, ymax=1.75, ymin=-0.73]
		\tkzDrawX[ noticks, label={ },above=5pt, line width=2pt ] %noticks, 
		\tkzClip
		\tkzDefPoint(-3.0,0){A1}
		\tkzDefPoint(2.0,0){A2}
		\tkzDefPoint(1.0,0){B1}
		\tkzDefPoint(5.0,0){B2}
		\draw    (A1) to[out=30,in=150]  node[pos=0.5, fill=white ] {$A$} (A2)  ;
		\draw    (B1) to[out=30,in=150]  node[pos=0.5, fill=white ] {$B$} (B2)  ;
		\tkzLabelPoint[below](A1){$-3$}
		\tkzLabelPoint[below](A2){$2$}
		\tkzDrawPoint[shape=circle,fill=white,size=8](A1);
		\tkzDrawPoint[shape=circle,fill=white,size=8](A2);
		\tkzLabelPoint[below](B1){$1$}
		\tkzLabelPoint[below](B2){$5$}
		\tkzDrawPoint[shape=circle, size=8](B1);
		\tkzDrawPoint[shape=circle,  size=8](B2);
	\end{tikzpicture} 	
}

\vspace{-0.75em}
\noindent
\parbox{0.4\linewidth}{
	\noindent $A=\interval[open]{-3}{1}$ et $B=\interval[]{2}{5}$.
	
	$A\cap B=\emptyset$ \\
	$A\cup B=\interval[open]{-3}{1}\cup \interval[]{2}{5}$ 
}
\hfill 
\parbox{0.6\linewidth}{
	\begin{tikzpicture}[xscale=0.8]
		\tkzInit[xmin=-6, xmax=6.0, xstep=1, ymax=1.75, ymin=-0.73]
		\tkzDrawX[ noticks, label={ },above=5pt, line width=2pt ] %noticks, 
		\tkzClip
		\tkzDefPoint(-3.0,0){A1}
		\tkzDefPoint(1.0,0){A2}
		\tkzDefPoint(2.0,0){B1}
		\tkzDefPoint(5.0,0){B2}
		\draw    (A1) to[out=30,in=150]  node[pos=0.5, fill=white ] {$A$} (A2)  ;
		\draw    (B1) to[out=30,in=150]  node[pos=0.5, fill=white ] {$B$} (B2)  ;
		\tkzLabelPoint[below](A1){$-3$}
		\tkzLabelPoint[below](A2){$1$}
		\tkzDrawPoint[shape=circle,fill=white,size=8](A1);
		\tkzDrawPoint[shape=circle,fill=white,size=8](A2);
		\tkzLabelPoint[below](B1){$2$}
		\tkzLabelPoint[below](B2){$5$}
		\tkzDrawPoint[shape=circle, size=8](B1);
		\tkzDrawPoint[shape=circle,  size=8](B2);
	\end{tikzpicture} 	
}
\end{example}


\begin{exercise} Pour chaque cas déterminez les nsembles $A\cap B$ et $A\cup B$ \hfill \\ 
 \noindent
 \parbox{0.4\linewidth}{
 	\noindent $A=\interval[open right]{-10}{2}$ et $B=\interval[]{-5}{3}$.
 	
 	$A\cap B= $\\
 	$A\cup B= $  
 }
 \hfill 
 \parbox{0.6\linewidth}{
 	\begin{tikzpicture}[xscale=0.8]
 		\tkzInit[xmin=-6, xmax=6.0, xstep=1, ymax=1.75, ymin=-0.73]
 		\tkzDrawX[ noticks, label={ },above=5pt, line width=2pt ] %noticks, 
 		\tkzClip
% 		\tkzDefPoint(-3.0,0){A1}
% 		\tkzDefPoint(2.0,0){A2}
% 		\tkzDefPoint(1.0,0){B1}
% 		\tkzDefPoint(5.0,0){B2}
% 		\draw    (A1) to[out=30,in=150]  node[pos=0.5, fill=white ] {$A$} (A2)  ;
% 		\draw    (B1) to[out=30,in=150]  node[pos=0.5, fill=white ] {$B$} (B2)  ;
% 		\tkzLabelPoint[below](A1){$-3$}
% 		\tkzLabelPoint[below](A2){$2$}
% 		\tkzDrawPoint[shape=circle,fill=white,size=8](A1);
% 		\tkzDrawPoint[shape=circle,fill=white,size=8](A2);
% 		\tkzLabelPoint[below](B1){$1$}
% 		\tkzLabelPoint[below](B2){$5$}
% 		\tkzDrawPoint[shape=circle, size=8](B1);
% 		\tkzDrawPoint[shape=circle,  size=8](B2);
 	\end{tikzpicture} 	
 }
 
  \noindent
 \parbox{0.4\linewidth}{
 	\noindent $A=\interval[open]{-\infty}{2}$ et $B=\interval[open right]{0}{5}$.
 	
 	$A\cap B= $\\
 	$A\cup B= $  
 }
 \hfill 
 \parbox{0.6\linewidth}{
 	\begin{tikzpicture}[xscale=0.8]
 		\tkzInit[xmin=-6, xmax=6.0, xstep=1, ymax=1.75, ymin=-0.73]
 		\tkzDrawX[ noticks, label={ },above=5pt, line width=2pt ] %noticks, 
 		\tkzClip
 		% 		\tkzDefPoint(-3.0,0){A1}
 		% 		\tkzDefPoint(2.0,0){A2}
 		% 		\tkzDefPoint(1.0,0){B1}
 		% 		\tkzDefPoint(5.0,0){B2}
 		% 		\draw    (A1) to[out=30,in=150]  node[pos=0.5, fill=white ] {$A$} (A2)  ;
 		% 		\draw    (B1) to[out=30,in=150]  node[pos=0.5, fill=white ] {$B$} (B2)  ;
 		% 		\tkzLabelPoint[below](A1){$-3$}
 		% 		\tkzLabelPoint[below](A2){$2$}
 		% 		\tkzDrawPoint[shape=circle,fill=white,size=8](A1);
 		% 		\tkzDrawPoint[shape=circle,fill=white,size=8](A2);
 		% 		\tkzLabelPoint[below](B1){$1$}
 		% 		\tkzLabelPoint[below](B2){$5$}
 		% 		\tkzDrawPoint[shape=circle, size=8](B1);
 		% 		\tkzDrawPoint[shape=circle,  size=8](B2);
 	\end{tikzpicture} 	
 }
 
\noindent
\parbox{0.4\linewidth}{
	\noindent $A=\interval[open right]{3}{\infty}$ et $B=\interval[open]{-\infty}{6}$.
	
	$A\cap B= $\\
	$A\cup B= $  
}
\hfill 
\parbox{0.6\linewidth}{
	\begin{tikzpicture}[xscale=0.8]
		\tkzInit[xmin=-6, xmax=6.0, xstep=1, ymax=1.75, ymin=-0.73]
		\tkzDrawX[ noticks, label={ },above=5pt, line width=2pt ] %noticks, 
		\tkzClip
		% 		\tkzDefPoint(-3.0,0){A1}
		% 		\tkzDefPoint(2.0,0){A2}
		% 		\tkzDefPoint(1.0,0){B1}
		% 		\tkzDefPoint(5.0,0){B2}
		% 		\draw    (A1) to[out=30,in=150]  node[pos=0.5, fill=white ] {$A$} (A2)  ;
		% 		\draw    (B1) to[out=30,in=150]  node[pos=0.5, fill=white ] {$B$} (B2)  ;
		% 		\tkzLabelPoint[below](A1){$-3$}
		% 		\tkzLabelPoint[below](A2){$2$}
		% 		\tkzDrawPoint[shape=circle,fill=white,size=8](A1);
		% 		\tkzDrawPoint[shape=circle,fill=white,size=8](A2);
		% 		\tkzLabelPoint[below](B1){$1$}
		% 		\tkzLabelPoint[below](B2){$5$}
		% 		\tkzDrawPoint[shape=circle, size=8](B1);
		% 		\tkzDrawPoint[shape=circle,  size=8](B2);
	\end{tikzpicture} 	
}

\noindent
\parbox{0.4\linewidth}{
	\noindent $A=\interval[open]{-\infty}{-2}$ et $B=\interval[open]{-4}{3}$.
	
	$A\cap B= $\\
	$A\cup B= $  
}
\hfill 
\parbox{0.6\linewidth}{
	\begin{tikzpicture}[xscale=0.8]
		\tkzInit[xmin=-6, xmax=6.0, xstep=1, ymax=1.75, ymin=-0.73]
		\tkzDrawX[ noticks, label={ },above=5pt, line width=2pt ] %noticks, 
		\tkzClip
		% 		\tkzDefPoint(-3.0,0){A1}
		% 		\tkzDefPoint(2.0,0){A2}
		% 		\tkzDefPoint(1.0,0){B1}
		% 		\tkzDefPoint(5.0,0){B2}
		% 		\draw    (A1) to[out=30,in=150]  node[pos=0.5, fill=white ] {$A$} (A2)  ;
		% 		\draw    (B1) to[out=30,in=150]  node[pos=0.5, fill=white ] {$B$} (B2)  ;
		% 		\tkzLabelPoint[below](A1){$-3$}
		% 		\tkzLabelPoint[below](A2){$2$}
		% 		\tkzDrawPoint[shape=circle,fill=white,size=8](A1);
		% 		\tkzDrawPoint[shape=circle,fill=white,size=8](A2);
		% 		\tkzLabelPoint[below](B1){$1$}
		% 		\tkzLabelPoint[below](B2){$5$}
		% 		\tkzDrawPoint[shape=circle, size=8](B1);
		% 		\tkzDrawPoint[shape=circle,  size=8](B2);
	\end{tikzpicture} 	
}

\noindent
\parbox{0.4\linewidth}{
	\noindent $A=\interval[open left]{-4}{2}$ et $B=\interval[]{2}{5}$.
	
	$A\cap B= $\\
	$A\cup B= $  
}
\hfill 
\parbox{0.6\linewidth}{
	\begin{tikzpicture}[xscale=0.8]
		\tkzInit[xmin=-6, xmax=6.0, xstep=1, ymax=1.75, ymin=-0.73]
		\tkzDrawX[ noticks, label={ },above=5pt, line width=2pt ] %noticks, 
		\tkzClip
		% 		\tkzDefPoint(-3.0,0){A1}
		% 		\tkzDefPoint(2.0,0){A2}
		% 		\tkzDefPoint(1.0,0){B1}
		% 		\tkzDefPoint(5.0,0){B2}
		% 		\draw    (A1) to[out=30,in=150]  node[pos=0.5, fill=white ] {$A$} (A2)  ;
		% 		\draw    (B1) to[out=30,in=150]  node[pos=0.5, fill=white ] {$B$} (B2)  ;
		% 		\tkzLabelPoint[below](A1){$-3$}
		% 		\tkzLabelPoint[below](A2){$2$}
		% 		\tkzDrawPoint[shape=circle,fill=white,size=8](A1);
		% 		\tkzDrawPoint[shape=circle,fill=white,size=8](A2);
		% 		\tkzLabelPoint[below](B1){$1$}
		% 		\tkzLabelPoint[below](B2){$5$}
		% 		\tkzDrawPoint[shape=circle, size=8](B1);
		% 		\tkzDrawPoint[shape=circle,  size=8](B2);
	\end{tikzpicture} 	
}

\noindent
\parbox{0.4\linewidth}{
	\noindent $A=\interval[]{-4}{2}$ et $B=\interval[open left]{2}{5}$.
	
	$A\cap B= $\\
	$A\cup B= $  
}
\hfill 
\parbox{0.6\linewidth}{
	\begin{tikzpicture}[xscale=0.8]
		\tkzInit[xmin=-6, xmax=6.0, xstep=1, ymax=1.75, ymin=-0.73]
		\tkzDrawX[ noticks, label={ },above=5pt, line width=2pt ] %noticks, 
		\tkzClip
		% 		\tkzDefPoint(-3.0,0){A1}
		% 		\tkzDefPoint(2.0,0){A2}
		% 		\tkzDefPoint(1.0,0){B1}
		% 		\tkzDefPoint(5.0,0){B2}
		% 		\draw    (A1) to[out=30,in=150]  node[pos=0.5, fill=white ] {$A$} (A2)  ;
		% 		\draw    (B1) to[out=30,in=150]  node[pos=0.5, fill=white ] {$B$} (B2)  ;
		% 		\tkzLabelPoint[below](A1){$-3$}
		% 		\tkzLabelPoint[below](A2){$2$}
		% 		\tkzDrawPoint[shape=circle,fill=white,size=8](A1);
		% 		\tkzDrawPoint[shape=circle,fill=white,size=8](A2);
		% 		\tkzLabelPoint[below](B1){$1$}
		% 		\tkzLabelPoint[below](B2){$5$}
		% 		\tkzDrawPoint[shape=circle, size=8](B1);
		% 		\tkzDrawPoint[shape=circle,  size=8](B2);
	\end{tikzpicture} 	
} 
\end{exercise}
\begin{exercise}[entrainement] Déterminez les intersections ci-dessous.
	\vspace{-1em}
	\begin{multicols}{4}
	\begin{enumerate}[leftmargin=*,label={}]
		\item $\interval{2}{5}\cap \interval[open right]{3}{6}$\dotfill 
		\item $\interval[open left]{-\infty}{3}\cap \interval[]{-7}{10}$\dotfill
		\item $\interval{-5}{2}\cup \interval{0}{5}$\dotfill
		\item $\interval{-2}{0}\cap \interval[open right]{4}{5}$\dotfill
%		\item $\interval[open left]{-\infty}{0}\cup \interval[open right]{0}{+\infty}$\dotfill
%		\item $\interval[open left]{-\infty}{0}\cap \interval[open right]{0}{+\infty}$\dotfill
%		\item $\interval{-2}{1}\cap \interval[open right]{-5}{0}\cap \interval[open right]{-1}{2}$\dotfill 
	\end{enumerate} 
	\end{multicols}
\end{exercise} 
%----------------------------------------------------------------------------------------
%	COURS
%----------------------------------------------------------------------------------------

\newpage 
\loadgeometry{marginlayout}  
\pagestyle{scrheadings} % Print headers again  
\KOMAoptions{
	headwidth=textwithmarginpar,
	headsepline=true,
	headsepline= 0.5pt,
	footwidth=textwithmarginpar,
} 
\onehalfspacing
\section{Relation d'ordre et opération}
\begin{definition}
	$a$ et $b\in \mathbb{R}$. \\
	$a$ est supérieur à $b$ s.s.i. la différence $(a-b)$ est positive :
	\[
	a\geqslant b \iff (a-b)\geqslant 0
	\]  
\end{definition} 
\marginnote{
	\begin{tBox}[leftline=false]
		Comparer deux expressions $a$ et $b$ revient à étudier le signe de la différence.
	\end{tBox}
}[-1cm]

\begin{theorem}[Addition] 
	L'addition conserve l'ordre :
	\[
	(a\geqslant b )\quad \Rightarrow \quad (a+n\geqslant b+n)
	\] 
\end{theorem}
\begin{proof}\hfill  
\end{proof}
\vfill
\begin{theorem}[Multiplication]  
	La \textbf{multiplication par un nombre positif} conserve l'ordre.\\
	La multiplication par un nombre \textbf{négatif} inverse l'ordre.
	\begin{align*}
		a\geqslant b \quad &\text{et } p\geqslant 0 \quad \Rightarrow \quad  pa \geqslant pb \\
		&  \text{et } n\leqslant 0 \quad \Rightarrow \quad na \leqslant nb 
	\end{align*}
\end{theorem}
\begin{proof}\hfill  
\end{proof}
\vfill

%----------------------------------------------------------------------------------------
%	EXERCICES
%----------------------------------------------------------------------------------------

\newpage 
\loadgeometry{widelayout}  
\pagestyle{scrheadings} % Print headers again  
\KOMAoptions{
	headwidth=textwithmarginpar,
	headsepline=true,
	headsepline= 0.5pt,
	footwidth=textwithmarginpar,
} 
\onehalfspacing
\subsection{Exercices : propriétés des inégalités}


\begin{exercise}\hfill 
\begin{enumerate}[label={\arabic*)}, leftmargin=*]
\item Complétez par $<$ ou $>$, $\geqslant$ et $\leqslant$.\vspace{-1em}
	\begin{spacing}{2}
	\begin{multicols}{2}
		\begin{enumerate}[label={}, leftmargin=*]
			\item Si $x<y$ alors $x+2 \ldots y+2$ et $y-x \ldots 0$.
			\item Si $x>0$ et $y\ldots 0$ alors $\frac{x}{y}<0$.
			\item Si $x>0$ et $y\ldots 0$ alors $xy\geqslant 0$.
			\item Si $4x-3>2y-3$ alors $4x\ldots 2y$.
			\item Si $3x-1<3$ alors $4 \ldots 3x$.
			\item Si $a<b$ alors $a-b \ldots b-a$.
		\end{enumerate}
	\end{multicols}
	\end{spacing}\vspace{-1em}
\item Vrai ou faux ? Si faux donner un contre-exemple.\\
\noindent\QCM[VF,Alterne,Stretch=1.3,  Largeur=1.0cm]{%NomV=Indépendants,NomF={\small Non indépendants}, 
	{Pour tout $a\in\mathbb{R}$, si $a^2>0$ alors $a>0$.}&2,% 
	{Pour tout $a\in\mathbb{R}$, si $a^2>a$ alors $a>0$.}&2,% 
	{Pour tout $a\in\mathbb{R}$, si $a<1$ alors $a^2<a$.}&2,% 
	{Pour tout $a\in\mathbb{R}$, si $a<0$ alors $a^2>a$.}&2,%  
} 
\item Vrai ou faux ? Si faux donner un contre-exemple.\\
\noindent\QCM[VF,Alterne,Stretch=1.3,  Largeur=1.0cm]{%NomV=Indépendants,NomF={\small Non indépendants},  
	{Si $-4x>20$ alors $x>-5$.}&2,% 
	{Si $a<b$ alors $ac<bc$.}&2,% 
	{Si $a>b$ alors $ac^2>bc^2$.}&1,% 
	{Si $ac^2>bc^2$ alors $a>b$.}&1,%  
	{Si $a>0$, $b<0$ et $c<0$ alors $(a-b)c>0$}&2,% 
} 
\item On donne les nombres $a$, $b$ et $c$ sur la droite graduée ci-dessous. Complétez par $<$ ou $>$, $\geqslant$ et $\leqslant$ :
\vspace{-1em}
\begin{spacing}{2}
	\begin{multicols}{3}
		\begin{enumerate}[label={}, leftmargin=*]
			\item 
			\begin{tikzpicture}[xscale=2.5]
				\tkzInit[xmin=-0.99,xmax=1.0,ymin=-0.5,ymax=0.5,xstep=1]
				\tkzDrawX[left space=0.05,right space=0.05,noticks,label={}];%tickwd=2pt,
				\tikzset{arr/.style={postaction=decorate,	decoration={markings,mark=at position 1 with {\arrow[thick]{#1}}}}}
				\foreach \x in {0,0.2,...,1.9}
				{\draw (-1.0+\x,-0.15)--(-1.0+\x,0.15);} 
				\tkzDefPoint(0,0){O}
				\tkzLabelPoint[color = black,below=5pt,inner sep = 5pt,font=\small](O){$0$} 
				\tkzDefPoint(0.63,0){A}
				\tkzLabelPoint[color = black,below=5pt,inner sep = 5pt,font=\small](A){$a$} 
				\tkzDefPoint(-0.2,0){B}
				\tkzLabelPoint[color = black,below=5pt,inner sep = 5pt,font=\small](B){$b$} 
				\tkzDefPoint(-0.82,0){C}
				\tkzLabelPoint[color = black,below=5pt,inner sep = 5pt,font=\small](C){$c$} 
			\end{tikzpicture} 
			\item $a+c\ldots\ldots  b+c$
			\item $\frac{a}{b}\ldots\ldots 0$
			\item $ab\ldots\ldots 0$
			\item $\frac{c}{b}\ldots\ldots 1$
		\end{enumerate}
	\end{multicols}
\end{spacing}\vspace{-1em}
\end{enumerate} 
\end{exercise} 
\begin{exercise}Complétez pour avoir des inégalités équivalentes.
	
\noindent
\parbox{0.37\linewidth}{
\begin{spacing}{2}
$\begin{WithArrows}
	-3x& > -3y \Arrow{$\times \ldots\ldots$}\\
\iff	-12x& \ldots\ldots -12y
\end{WithArrows}$
\end{spacing}
}
\parbox{0.31\linewidth}{
\begin{spacing}{2}
	$\begin{WithArrows}
		b& <c \Arrow{$a<0$}\\
\iff	ab& \ldots\ldots ac
	\end{WithArrows}$
\end{spacing}
} 
\parbox{0.31\linewidth}{
	\begin{spacing}{2}
		$\begin{WithArrows}
			a& <2b \Arrow{$-c$}\\
			\iff \qquad & \ldots \ldots \qquad
		\end{WithArrows}$
	\end{spacing}
}	


\noindent
\parbox{0.4\linewidth}{
	\begin{spacing}{2}
		$\begin{WithArrows}
			x-2y& >x \Arrow{$-x$}\\
			\iff	-2y& \ldots\ldots 0\Arrow{$\times\ldots\ldots$} \\
			\iff y & \ldots\ldots 0
		\end{WithArrows}$
	\end{spacing}
}	
\parbox{0.4\linewidth}{
	\begin{spacing}{2}
		$\begin{WithArrows}
			x & <y \Arrow{$\times \tfrac{-2}{3}$}\\
			\iff	-\frac{2}{3}x& \ldots\ldots -\frac{2}{3}y\Arrow{$+1$} \\
			\iff -\frac{2}{3}x+1 & \ldots\ldots -\frac{2}{3}y+1
		\end{WithArrows}$
	\end{spacing}
}	
\end{exercise}
%----------------------------------------------------------------------------------------
%	EXERCICES
%----------------------------------------------------------------------------------------

\newpage 
\loadgeometry{widelayout}  
\pagestyle{scrheadings} % Print headers again  
\KOMAoptions{
	headwidth=textwithmarginpar,
	headsepline=true,
	headsepline= 0.5pt,
	footwidth=textwithmarginpar,
} 
\onehalfspacing
\subsection{Exercices : résolution d'inéquations du premier degré}
\begin{tBox}
On ne change pas les solutions d'une inéquation si :
	\begin{enumerate}[label=\textbullet, leftmargin=*]
		\item on \textbf{développe, factorise, réduit} un des membres de l'inéquation
		\item on \textbf{ajoute}  \textbf{une même expression}  aux deux membres de 
		\item on \textbf{multiplie} les 2 membres de l'inéquation par \textbf{une même expression positive non nulle}
		\item  on \textbf{multiplie} les 2 membres de l'inéquation par \textbf{une même expression négative non nul(le)} à condition de  \textbf{changer le sens du signe de l'inéquation}
	\end{enumerate}  
\end{tBox}  
\begin{example}  Résoudre dans $\mathbb{R}$ les inéquations suivantes d'inconnue $x$ :\hfill 
	
	\noindent 
	\parbox{0.31\linewidth}{
		\begin{spacing}{2}
			$\begin{WithArrows}
				3x+4&>10 \Arrow{$-4$}\\
				\iff 3x\qquad& \ldots   6 \Arrow{$\times \tfrac{1}{3}$}\\
				\iff\frac{3}{3} x\qquad& \ldots  \tfrac{6}{3} \\
				\iff x & \ldots\ldots  
			\end{WithArrows}$\\
		$\mathscr{S}=$
		\end{spacing}
	}	
%	
%	
%	\parbox{0.31\linewidth}{	
%		$\begin{WithArrows}
%			3x+4&>10 \\
%			&  \rule{0pt}{1.3em}   \\
%			&  \rule{0pt}{1.3em}   \\
%			&  \rule{0pt}{1.3em}    \\
%			&  \rule{0pt}{1.3em}  
%		\end{WithArrows}$ 	
%	}
\hfill 
	\parbox{0.31\linewidth}{
	\begin{spacing}{2}
		$\begin{WithArrows}
		-2x-8&>10 \Arrow{$+8$}\\
			\iff \qquad & \ldots   \qquad \Arrow{$\times \tfrac{1}{-2}$}\\
			\iff\frac{-2}{-2} x& \ldots  \qquad\qquad \\
			\iff x & \ldots   \qquad\qquad
		\end{WithArrows}$\\
		$\mathscr{S}=$
	\end{spacing}
}	
%
%\parbox{0.31\linewidth}{	
%		$\begin{WithArrows}
%			-2x-8&>10 \\
%			&  \rule{0pt}{1.3em}   \\
%			&  \rule{0pt}{1.3em}   \\
%			&  \rule{0pt}{1.3em}  \\
%			&  \rule{0pt}{1.3em}  
%		\end{WithArrows}$ 
%	} 
\hfill
	\parbox{0.31\linewidth}{
	\begin{spacing}{2}
		$\begin{WithArrows}
			-3x-8&\leqslant 10 \Arrow{$+\ldots$}\\
			\iff \quad\qquad & \ldots   \qquad \Arrow{$\times \ldots$}\\
			\iff  \quad\qquad & \ldots  \qquad\qquad \\
			\iff \quad\qquad & \ldots   \qquad\qquad
		\end{WithArrows}$\\
		$\mathscr{S}=$
	\end{spacing}
}	
% \parbox{0.31\linewidth}{	
%		$\begin{WithArrows}
%			2+5x&\leqslant -13 \\
%			&  \rule{0pt}{1.3em}   \\
%			&  \rule{0pt}{1.3em}   \\
%			&  \rule{0pt}{1.3em}   \\
%			&  \rule{0pt}{1.3em}  
%		\end{WithArrows}$ 
%	}    
\end{example}
\begin{exercise}[variables d'un seul côté]\label{chapter04:exercice:inequation01} Même consignes
	\vspace{-1em}
	\begin{multicols}{3}
		\begin{enumerate}[label={$(I_{\arabic*})$\rule{0pt}{1.7em}}, leftmargin=*]
			\item $x+1<9$
			\item $x-4>3$
%			\item $x+10\leqslant 12$
			\item $-6x\geqslant 30$
			\item $-x<8$
			\item $7<2x-11$
			\item $-8x-5>0$
			\item $42x>0$
			\item $14-6x\geqslant -10$
			\item $\tfrac{3}{2}x-1>4$
%			\item $-5<\frac{2}{5}x$
%			\item $6-\frac{x}{5}>-14$ 
		\end{enumerate}
	\end{multicols}  
\end{exercise}
 
\begin{exercise}[variable dans les deux membres] \label{chapter04:exercice:inequation02}  Mêmes consignes \hfill 
	\vspace{-1em}
	\begin{multicols}{3}
		\begin{enumerate}[label={$(I_{\arabic*})$\rule{0pt}{1.7em}}, leftmargin=*] 
			\item $3x>2x+1$ 
			\item $12x\leqslant 8x+128$
			\item $x+5<10x$  
			\item $3x+1\geqslant 3x+7$
			\item $3(x+1)-30<x+15$
			\item $2x-10<7x+5$
%			\item $6x-6\geqslant 3x+2$
			\item $5x+9\geqslant 5x+2$
			\item $1-7x\leqslant 7+x$
			\item $5x-5>-9x+3$
		\end{enumerate}
	\end{multicols}  
\end{exercise}
%\item $x+7\geqslant -5x-1$ 
%\item $2x<20$
%\item $2x<-20$
%\item $-2x\geqslant 20$
%\item $x-5<0$
%\item $-x+4\geqslant 0$
%\item $5<\frac{2}{3}x$
%\item $\tfrac{3}{4}x>24$ 
\begin{example} Expliquer les erreurs dans les résolutions suivantes  \hfill 
	
	
	\noindent\parbox{0.45\linewidth}{$\begin{WithArrows}
			\frac{1}{x}&\geqslant -1 \Arrow{On multiplie par $x$}\\
			\iff\quad	\frac{x}{x}&\geqslant -x\Arrow{On simplifie}\\
			\iff\quad	1&\geqslant -x\Arrow{On multiplie par $-1$}\\
			\iff\quad	-1&\leqslant x \\
		 \mathscr{S}&=\interval[open right, scaled]{-1}{+\infty}
		\end{WithArrows}$ 
	}\hfill
\parbox{0.45\linewidth}{$\begin{WithArrows}
		x^2&> 5x \Arrow{On divise par $x$}\\
		\iff\quad	\frac{x^2}{x}&>\frac{5x}{x}\Arrow{On simplifie}\\
		\iff\quad	x& > 5\\
		 \mathscr{S}&=\interval[open, scaled]{5}{+\infty}
	\end{WithArrows}$ 
} 
\end{example} 
\begin{exercise}\label{chapter04:exercice:inequation02bonus}     Mêmes consignes \hfill 
	\vspace{-1em}
	\begin{multicols}{2}
		\begin{enumerate}[label={$(I_{\arabic*})$\rule{0pt}{1.7em}}, leftmargin=*] 
			\item $5(x-1)>4(2x-1)$ 
			\item $\frac{3x-1}{4}-1\geqslant 0$
			\item $2(1-x)+3(-1-x)\geqslant 6-(3x+2)$ 
			\item $5-4[3(x-2)-1]<57$ 
		\end{enumerate}
	\end{multicols}  
\end{exercise}
\newpage 
\begin{example}[Encadrements]  Résoudre dans $\mathbb{R}$ les inéquations suivantes d'inconnue $x$ :\hfill 
	
	\noindent \parbox{0.31\linewidth}{	
		$\begin{WithArrows}
			-1&<x+2 < 5 \\
			&  \rule{0pt}{1.3em}   \\
			&  \rule{0pt}{1.3em}   \\
			&  \rule{0pt}{1.3em}    \\
			&  \rule{0pt}{1.3em}    
		\end{WithArrows}$ 	
	}\hfill \parbox{0.31\linewidth}{	
		$\begin{WithArrows}
			-4&\geqslant -2x \geqslant - 10 \\
			&  \rule{0pt}{1.3em}   \\
			&  \rule{0pt}{1.3em}   \\
			&  \rule{0pt}{1.3em}  \\
			&  \rule{0pt}{1.3em} 
		\end{WithArrows}$ 
	} \hfill \parbox{0.31\linewidth}{	
		$\begin{WithArrows}
			-5&\leqslant -3x+7 <15 \\
			&  \rule{0pt}{1.3em}   \\
			&  \rule{0pt}{1.3em}   \\
			&  \rule{0pt}{1.3em}   \\
			&  \rule{0pt}{1.3em}   
		\end{WithArrows}$ 
	}    
\end{example}
\vfill 
\begin{exercise}  \label{chapter04:exercice:inequation03}  Mêmes consignes \hfill 
	\vspace{-1em} 
	\begin{multicols}{3}
		\begin{enumerate}[label={$(I_{\arabic*})$\rule{0pt}{1.4em}}, leftmargin=*] 
			\item $-3<x-4<7$ 
			\item $4<5x-4\leqslant 5$
			\item $-6\leqslant 3+x<4$  
			\item $2\leqslant 2x<10$
			\item $-1\leqslant -x<3$
			\item $-3\leqslant 1-x<4$ 
			\item $-3\leqslant 2x-1<1$
			\item $8<-2+3x<16$
			\item $4<2x-1\leqslant 10$
		\end{enumerate}
	\end{multicols}  
\end{exercise}
\begin{example}[Disjonctions et inéquations simultanées]  Résoudre dans $\mathbb{R}$ les inéquations suivantes   :\hfill 
	
	\noindent \parbox{0.31\linewidth}{	
		$\begin{WithArrows}
			3&<x  \text{ et } x<7 \\
			&  \rule{0pt}{1.3em}   \\
			&  \rule{0pt}{1.3em}   \\
			&  \rule{0pt}{1.3em}    \\
			&  \rule{0pt}{1.3em}     
		\end{WithArrows}$ 	
	}\hfill \parbox{0.31\linewidth}{	
		$\begin{WithArrows}
			x&>4  \text{ ou } x<3 \\
			&  \rule{0pt}{1.3em}   \\
			&  \rule{0pt}{1.3em}   \\
			&  \rule{0pt}{1.3em}  \\
			&  \rule{0pt}{1.3em}   
		\end{WithArrows}$ 
	}   \hfill \parbox{0.31\linewidth}{	
	$\begin{WithArrows}
		x&>1  \text{ ou } x=-1 \\
		&  \rule{0pt}{1.3em}   \\
		&  \rule{0pt}{1.3em}   \\
		&  \rule{0pt}{1.3em}  \\
		&  \rule{0pt}{1.3em}   
	\end{WithArrows}$ 
}     
\end{example}
\begin{exercise}  \label{chapter04:exercice:inequation04}  Mêmes consignes \hfill
	\vspace{-1em} 
	\begin{multicols}{3}
\begin{enumerate}[label={$(I_{\arabic*})$\rule{0pt}{1.4em}},leftmargin=*]
	\item  $x\leqslant -3$ ou $x\geqslant 1$
	\item $x< 4$ ou $x> 8$
	\item  $x< 9$ et $x< -3$
%	\item  $x< 9$ ou $x< -3$
	\item  $x\leqslant 1$ et $x> 5$
	\item  $x< -3$ ou $x= 5$
%	\item  $x+2\leqslant -1$ et $x+2\geqslant 3$
%	\item  $-2x>10$ ou $4x>16$
%	\item  $-3x+1>10$ ou $2x+1>11$
%	\item $15>4x-1$ ou $1<4x-15$  
	\item $x+5\leqslant -4$ ou $x+5\geqslant 4$
	\item $-2x>10$ ou $4x>16$
%	\item $x-3>6$ ou $x+12<9$
	\item $15>4x-1$ ou $1<4x-15$ 
	\item $3x+1\leqslant 4$ et $2x-3>7$
\end{enumerate}
\end{multicols}  
\end{exercise}
%\begin{cBox}
% Quel est l'ensemble solution de \og  $15>4x-1$ ou $1<4x-15$ \fg{}, inconnue $x$ ?
%	\end{cBox}
\begin{example}[Valeur absolue]  Résoudre dans $\mathbb{R}$ les inéquations suivantes d'inconnue $x$ :\hfill 
	
	\noindent \parbox{0.45\linewidth}{	
		$\begin{WithArrows}
			\abs{x+3}&<6   \\
			&  \rule{0pt}{1.3em}   \\
			&  \rule{0pt}{1.3em}   \\
			&  \rule{0pt}{1.3em}      
		\end{WithArrows}$ 	
	}\hfill \parbox{0.45\linewidth}{	
		$\begin{WithArrows}
			\abs{2x}&\geqslant 10 \\
			&  \rule{0pt}{1.3em}   \\
			&  \rule{0pt}{1.3em}   \\
			&  \rule{0pt}{1.3em}     
		\end{WithArrows}$ 
	}   
\end{example}
\vfill 
\begin{exercise}  \label{chapter04:exercice:inequation05}  Mêmes consignes \hfill 
	\vspace{-1em}
	\begin{multicols}{3}
		\begin{enumerate}[label={$(I_{\arabic*})$\rule{0pt}{1.4em}},leftmargin=*]
			\item $\abs{x+6}>7$
			\item $\abs{x+3}<4$
			\item $\abs{6x}>12$
			\item $\abs{1+2x}\geqslant 23$
			\item $\abs{2x-5}<7$
			\item $\abs{\tfrac{1}{4}x}>12$ 
		\end{enumerate}
	\end{multicols}  
\end{exercise}
%----------------------------------------------------------------------------------------
%	solutions
%----------------------------------------------------------------------------------------

\newpage 
\begin{problem} Un artisant doit fabriquer 176 jouets le mois de juillet. Durant les 10 premiers jours il en fait 4 par jours. Après l'achat de nouveaux outils, il termine la commande 3 jours avant la date prévue. On désigne par $x$ le nombre moyen de jouets produit chaque jour après le 10\ieme jour. 
	\begin{enumerate}[leftmargin=*,label=\alph*)]
		\item Écrire une inéquation vérifiée par $x$.
		\item Déterminez la valeur minimale de $x$.
	\end{enumerate} 
	
\end{problem}
\begin{proof} [solution de l'exercice  \ref{chapter04:exercice:inequation01}]\hfill 
	\vspace{-0.75em}
\begin{multicols}{3} 
\begin{pycode}
from re import sub
import sympy as sp
from sympy.core.sympify import kernS

x, y, z, t , a, b = sp.symbols('x y z t a b')
k, m, n = sp.symbols('k m n', integer=True)
f, g, h = sp.symbols('f g h', cls=sp.Function)

# Create a list of functions to include in the table
funcs = [
'x+1<9',
'x-4>3', 
'-6*x>= 30',
'-x<8',
'7<2*x-11',
'-8*x-5>0',
'42*x>0',
'14-6*x>= -10',
'3/2*x-1>4'
]
n=1
print(r'\begin{enumerate}[label={$(I_{\arabic*})$}] ')
for func in funcs:
	print('\item ')
	print( '$\mathscr{S}  = ' + sp.latex(sp.solveset(func,   x,   domain=sp.S.Reals)).replace("(", "]").replace(")", "[")   + '$'     )
print(r'\end{enumerate}  ')
\end{pycode}
\end{multicols}	
\scriptsize
\end{proof} 
\vspace{-1em}
\begin{proof} [solution de l'exercice  \ref{chapter04:exercice:inequation02}]\hfill
	\vspace{-0.75em}
\begin{multicols}{3} 
\begin{pycode}
from re import sub
import sympy as sp
from sympy.core.sympify import kernS

x, y, z, t , a, b = sp.symbols('x y z t a b')
k, m, n = sp.symbols('k m n', integer=True)
f, g, h = sp.symbols('f g h', cls=sp.Function)

# Create a list of functions to include in the table
funcs = [
'3*x>2*x+1', 
'12*x<= 8*x+128',
'x+5<10*x',  
'3*x+1>= 3*x+7',
'3*(x+1)-30<x+15',
'2*x-10<7*x+5', 
'5*x+9>= 5*x+2',
'1-7*x<= 7+x',
'5*x-5>-9*x+3',
]
n=1
print(r'\begin{enumerate}[label={$(I_{\arabic*})$}] ')
for func in funcs:
	print('\item ')
	print( '$\mathscr{S}  = ' + sp.latex(sp.solveset(func,   x,   domain=sp.S.Reals)).replace("(", "]").replace(")", "[")   + '$'     )
print(r'\end{enumerate}  ')
\end{pycode}
\end{multicols}	
\scriptsize
\end{proof}  
\vspace{-1em}
\begin{proof} [solution de l'exercice  \ref{chapter04:exercice:inequation02bonus}]\hfill 
	\vspace{-0.75em}
\begin{multicols}{4} 
\begin{pycode}
from re import sub
import sympy as sp
from sympy.core.sympify import kernS

x, y, z, t , a, b = sp.symbols('x y z t a b')
k, m, n = sp.symbols('k m n', integer=True)
f, g, h = sp.symbols('f g h', cls=sp.Function)

# Create a list of functions to include in the table
funcs = [
'5*(x-1)>4*(2*x-1)', 
'(3*x-1)/4>=1', 
'2*(1-x)+3*(-1-x)>=6- (3*x+2)', 
'5-4*(3*(x-2)-1)<57',  
]
n=1
print(r'\begin{enumerate}[label={$(I_{\arabic*})$}] ')
for func in funcs:
	print('\item ')
	print( '$\mathscr{S}  = ' + sp.latex(sp.solveset(func,   x,   domain=sp.S.Reals)).replace("(", "]").replace(")", "[")   + '$'     )
print(r'\end{enumerate}  ')
\end{pycode}
\end{multicols}	
\scriptsize
\end{proof} 
\vspace{-1em}
\begin{proof} [solution de l'exercice  \ref{chapter04:exercice:inequation03}]\hfill
	\vspace{-0.75em}
\begin{multicols}{3} 
\begin{pycode} 
from re import sub
import sympy as sp
from sympy.core.sympify import kernS

x, y, z, t , a, b = sp.symbols('x y z t a b')
k, m, n = sp.symbols('k m n', integer=True)
f, g, h = sp.symbols('f g h', cls=sp.Function)

# Create a list of functions to include in the table
funcs = [ 
['-3<x-4', 'x-4<7', 'and'],
['4<5*x-4', '5*x-4<= 5', 'and'],
['-6<= 3+x', '3+x<4', 'and'],
['2<= 2*x', '2*x<10', 'and'],
['-1<= -x', '-x<3', 'and'],
['-3<= 1-x', '1-x<4', 'and'],
['-3<= 2*x-1', '2*x-1<1', 'and'],
['8<-2+3*x', '-2+3*x<16', 'and'] ,
['4<2*x-1', '2*x-1<= 10', 'and'], 
]
n=1
print(r'\begin{enumerate}[label={$(I_{\arabic*})$}] ')
for func in funcs:
	print('\item ')
	A = sp.solveset(func[0],   x,   domain=sp.S.Reals)
	B = sp.solveset(func[1],   x,   domain=sp.S.Reals)
	AetB = sp.Intersection(A,B)
	print( '$\mathscr{S} =' + sp.latex(AetB).replace("(", "]").replace(")", "[")   + '$'     )
print(r'\end{enumerate} ')
\end{pycode}
\end{multicols}	
\scriptsize
\end{proof} 
\vspace{-1em}
\begin{proof}[solution de l'exercice  \ref{chapter04:exercice:inequation04}] \hfill
	\vspace{-0.75em}  
\begin{multicols}{3} 
\begin{pycode} 
from re import sub
import sympy as sp
from sympy.core.sympify import kernS

x, y, z, t , a, b = sp.symbols('x y z t a b')
k, m, n = sp.symbols('k m n', integer=True)
f, g, h = sp.symbols('f g h', cls=sp.Function)

# Create a list of functions to include in the table
funcs = [
['x<=-3', 'x>=1', 'or'],
['x< 4', 'x> 8', 'or'],
['x< 9', 'x< -3', 'and'],
['x<=1', 'x>5', 'and'],
['x<-3', '(x-5)**2<=0', 'or'],
['x+5<-4', 'x+5>=4', 'or'],
['-2*x>10', '4*x>16', 'or'],
['15>4*x-1', '1<4*x-15', 'or'] ,
['3*x+1<=4', '2*x-3>7', 'and'], 
]
n=1
print(r'\begin{enumerate}[label={$(I_{\arabic*})$}] ')
for func in funcs:
	print('\item ')
	A = sp.solveset(func[0],   x,   domain=sp.S.Reals)
	B = sp.solveset(func[1],   x,   domain=sp.S.Reals)
	AetB = sp.Intersection(A,B)
	AouB = sp.Union(A,B)
	if func[2]=='or' :
		print( '$\mathscr{S} =' + sp.latex(AouB).replace("(", "]").replace(")", "[")   + '$'     )
	elif func[2]=='and' :
		print( '$\mathscr{S} =' + sp.latex(AetB).replace("(", "]").replace(")", "[")   + '$'     )
print(r'\end{enumerate} ')
\end{pycode}
\end{multicols}
\scriptsize 
\end{proof} 
\vspace{-1em}
\begin{proof} [solution de l'exercice  \ref{chapter04:exercice:inequation05}]\hfill 
	\vspace{-0.5em}
\begin{multicols}{3} 
\begin{pycode}
from re import sub
import sympy as sp
from sympy.core.sympify import kernS
from sympy.solvers.inequalities import solve_univariate_inequality

x, y, z, t , a, b = sp.symbols('x y z t a b', real=True)
k, m, n = sp.symbols('k m n', integer=True)
f, g, h = sp.symbols('f g h', cls=sp.Function)

# Create a list of functions to include in the table
funcs = [
abs(x+6)>7, 
abs(x+3)<4, 
abs(6*x)>12, 
abs(1+2*x)>=23, 
abs(2*x-5)<7, 
abs(x/4)>12,  
]
n=1
print(r'\begin{enumerate}[label={$(I_{\arabic*})$}] ')
for func in funcs:
	print('\item ')
	print( '$\mathscr{S} = ' + sp.latex(solve_univariate_inequality(func,   x,   domain=sp.S.Reals, relational=False)).replace("(", "]").replace(")", "[")   + '$'     )
print(r'\end{enumerate}  ')
\end{pycode}
\end{multicols}	
\scriptsize
\end{proof} 
% \begin{exercise} \hfill
%	
%	Walid doit louer une voiture pour ses week-end. L'agence locale lui propose 3 formules :
%	\begin{itemize}
%		\item Formule \og Week-end \fg{} : \EUR{150} quel que soit le nombre de kilomètres parcourus.
%		\item Formule \og À l'unité \fg{} : \EUR{1.5} par \si{\kilo\meter} parcouru.
%		\item Formule \og Plus \fg{} : \EUR{100} plus \EUR{0.2} par \si{\kilo\meter} parcouru.
%	\end{itemize}
%	Determine,  pour chaque formule, le nombre de kilométrage qui la rendent la moins chère.
%\end{exercise}
%----------------------------------------------------------------------------------------
%	COURS
%----------------------------------------------------------------------------------------

\newpage 
\loadgeometry{marginlayout}  
\pagestyle{scrheadings} % Print headers again  
\KOMAoptions{
	headwidth=textwithmarginpar,
	headsepline=true,
	headsepline= 0.5pt,
	footwidth=textwithmarginpar,
} 
\onehalfspacing
\section{Bilan valeur absolue}

 
\begin{property} Pour $a$ et $r\in\mathbb{R}$. 
	\setlength{\abovedisplayskip}{3pt}
	\setlength{\belowdisplayskip}{3pt}
$$\abs{x-a}\leqslant r \qquad \text{signifie}\qquad x\in\interval{r-a}{r+a}$$
Intervalle centré en $a$
\end{property}
\begin{marginfigure}  
	\begin{tikzpicture}[scale=0.50] 
		\tkzInit[xmin=-4.5, xmax=4.50, xstep=1, ymax=0.75, ymin=-0.75]
		\tkzDrawX[  noticks, label={\scriptsize$x$},above=5pt ] %noticks, 
		%	  \tkzClip
		\tkzHTick[mark=text,mark options={color=black}, text mark={\large |}]{4.01}
		\tkzHTick[mark=text,mark options={color=black}, text mark={ {\large |}}]{-2.01}
		\tkzHTick[mark=text,mark options={color=black}, text mark={ {\large |}}]{0.01}
		\tkzDefPoint(0,0){O}
		\tkzDefPoint(-2,0){A}
		\tkzDefPoint(1.85,0){B}
		\tkzDrawSegment[dim={\scriptsize $\abs{x-a}$,+10pt,}](A,B)
		%				\tkzDrawSegments[  line width=3 pt, color=black](A,B)
		\tkzLabelPoint[below=2pt](A){\scriptsize$a=-2$}
		\tkzLabelPoint[below=2pt](B){\scriptsize$x$}
		\tkzLabelPoint[below=2pt](O){\scriptsize$0$}
	\end{tikzpicture} 
	\caption{L'écart entre $x$ et $a=-2$.}  
\end{marginfigure}
\begin{example} Il est plus pratique de retrouver l'intervalle en résolvant l'inéquations proprement comme vu en exercice.\\
	\noindent
\parbox{0.48\linewidth}{
$\begin{WithArrows}
	\abs{x+2}&\leqslant 0.1\\
	-0.1\leqslant x+2& \leqslant 0.1\Arrow{$-2$}\\
	-2-0.1\leqslant x&\leqslant-2+0.1\\
	-2.1 \leqslant x&\leqslant -1.9
	\end{WithArrows}
$
}
\hfill 
\parbox{0.48\linewidth}{
	$\begin{WithArrows}
		\abs{x-3}&\leqslant 0.1\\
		-0.1\leqslant x-3& \leqslant 0.1\Arrow{$+3$}\\
		3-0.1\leqslant x&\leqslant 3+0.1\\
		2.9 \leqslant x&\leqslant 3.1
	\end{WithArrows}
	$
}	
	\end{example}
%----------------------------------------------------------------------------------------
%	 
%----------------------------------------------------------------------------------------
\end{document}